\documentclass[11pt]{article}

\usepackage[margin=1in]{geometry}
\usepackage{booktabs}
\usepackage{graphicx}
\usepackage{hyperref}
\usepackage{xcolor}
\usepackage{array}
\usepackage{amsmath}

\hypersetup{
    colorlinks=true,
    linkcolor=blue,
    urlcolor=blue,
    citecolor=blue
}

\newcommand{\code}[1]{\texttt{#1}}

\title{\textbf{Reproducing HLS-Eval and Analog System Benchmarks\\
from Cong Hao's Group}\\[0.4em]
\large Results, what broke, and a proposed first project}
\author{Michael Incorvaia}
\date{August 14, 2026}

\begin{document}

\maketitle

\begin{abstract}
This note reports two reproductions of work co-authored by Cong (Callie)
Hao, intended as a concrete reading of the lab rather than a survey.
First, I reproduced the zero-shot HLS generation baseline from
\emph{HLS-Eval} (Abi-Karam \& Hao, ICLAD 2025) on 57 designs without
Vitis HLS, then re-ran the same harness on three post-paper models.
DeepSeek-V3 lands at 46.3\% run@1 versus the paper's 63.3\%, but
63.2\% versus 65.9\% at run@5; the pass@1 gap is largely a fixed
\code{max\_tokens=8192} cap. Under that same cap, Qwen3-Coder reaches
67.4\% run@1. Second, I treated Hasler \& Hao's TRETS 2024 analog
benchmark paper and the Ige/Hasler/Hao JLPEA 2023 analog-HLS paper as
one project: re-derived the Fig.~3/11 scaling anchors and the Fig.~17--18
Case~II CAB arithmetic (128 WTA CABs / 98 SoC CABs $= 1.31$ chips ---
the ``barely more than one SoC FPAA, fits in two'' claim holds), and
executed TRETS Fig.~16's acoustic program as a digital macromodel.
On a 4-word Speech Commands subset the VMM+WTA classifier reaches 70\%
held-out accuracy (chance 25\%). The AnalogBenchmarks artifact URL cited
in TRETS footnote~1 was already 404 on the Wayback Machine in November
2023. I did not reproduce SmartPAF or CRYPTONITE; both are discussed as
future work, as is AMS-Eval --- an HLS-Eval-style LLM harness on the
Fig.~15 analog block library, which is the project I would want to start
in the group.
\end{abstract}

\tableofcontents
\vspace{0.8em}

\section{What this note is}

I read several of your papers with the aim of reproducing what can be
reproduced from outside the lab, without proprietary EDA or an FPAA
board, and saying plainly what cannot. Two threads are finished:

\begin{enumerate}
\item \textbf{HLS-Eval}~\cite{hlseval} --- a measurement paper. There is
a pass@$k$ table to match. Done 2026-08-12.
\item \textbf{TRETS 2024 analog benchmarks}~\cite{trets} folded with
\textbf{JLPEA 2023 analog HLS}~\cite{jlpea} --- a benchmark-definition
paper plus its predecessor. There is no pass@$k$ table. The reproducible
claims are arithmetic (scaling curves, CAB costs) and making Fig.~16
executable. Done 2026-08-14.
\end{enumerate}

A third thread is a \emph{proposal}, not a result: \textbf{AMS-Eval},
the HLS-Eval metric ladder applied to analog block programs. SmartPAF
and CRYPTONITE were read and are interesting; they were not run.
Nothing below is a silicon measurement, and synthesis (Vitis / FPAA
place-and-route) is out of scope throughout.

%----------------------------------------------------------------------
\section{HLS-Eval: zero-shot generation, one year later}
\label{sec:hls}

\subsection{What was reproduced}

The task is the paper's: given a natural-language kernel description, a
header, and a testbench, generate C++. I score three of four staged
metrics --- parseable, compilable, runnable --- at
$n_{\text{samples}}=5$, $T=0.7$. Synthesizable is skipped, not failed:
Vitis HLS is not installed. Every synthesis cell is recorded as
skipped.

\textbf{Scope.} 57 designs from PolyBench, MachSuite, and C2HLSC, the
three suites that compile and run at 100\% under stock \code{clang++}
with no Xilinx software. (A prior audit counted 60 such designs; the
sweep used 57. That 3-design gap was not chased.) CHStone and Rosetta
are excluded for missing testbenches and Xilinx-header reasons
documented in \code{FINDINGS.md}.

\textbf{Provenance.} Upstream \code{hls-eval} @ \code{adea9ff}
(2026-07-30), evaluated 2026-08-12, WSL2 Ubuntu, \code{clang++} 18.1.3,
OpenRouter. The paper used Together AI. OpenRouter load-balances across
providers and quantizations, so numeric precision is uncontrolled here
in a way it was not for the paper. 1{,}140 LLM calls; 2 provider errors,
both retried; 0 unrecovered.

The paper's other anchor, Qwen2.5-Coder-32B, was dropped before the
sweep: every real HLS-Eval prompt to its sole OpenRouter provider
(Cloudflare) returned \code{finish\_reason: error}, while a short
synthetic \code{curl} succeeded. DeepSeek-V3 is the sole paper anchor
in this run.

\subsection{Paper vs.\ this run}

\begin{table}[h]
\centering
\small
\begin{tabular}{@{}lcccccc@{}}
\toprule
\textbf{Model} & \textbf{parse@1} & \textbf{parse@5} & \textbf{compile@1} & \textbf{compile@5} & \textbf{run@1} & \textbf{run@5} \\
\midrule
DeepSeek V3 (paper, 94 dsn) & 100.0\% & 100.0\% & 94.1\% & 97.6\% & \textbf{63.3\%} & \textbf{65.9\%} \\
Qwen2.5-Coder-32B (paper) & 95.5\% & 97.6\% & 90.8\% & 92.9\% & 57.6\% & 58.8\% \\
\midrule
deepseek-chat-v3-0324 (this run, 57 dsn)
  & 63.2\% & 87.7\% & 61.4\% & 86.0\% & 46.3\% & \textbf{63.2\%} \\
\emph{95\% CI}
  & \scriptsize{[54.0, 72.3]}
  & \scriptsize{[79.0, 94.7]}
  & \scriptsize{[52.3, 70.2]}
  & \scriptsize{[77.2, 94.7]}
  & \scriptsize{[35.8, 56.8]}
  & \scriptsize{[50.9, 75.4]} \\
\bottomrule
\end{tabular}
\caption{Paper Table~III (Together AI, 94 designs, Vitis included) vs.\
this run (OpenRouter, 57 designs, no synthesis). Bootstrap CIs over
designs, 10{,}000 resamples.}
\label{tab:hls-repro}
\end{table}

The pass@1 gap is real; the pass@5 gap on runnable nearly vanishes
(63.2\% vs.\ 65.9\%). A model that recovers by the fifth sample but
fails the first is a budget or formatting problem, not a capability
gap. \code{max\_tokens=8192} is not in the paper's stated methodology
and terminated 21\% of this run's calls with
\code{finish\_reason=length}. Pass@1 numbers here are a \textbf{lower
bound} on model capability.

Three confounds, stated plainly: (i)~57 designs vs.\ 94; (ii)~OpenRouter
vs.\ Together AI; (iii)~the token cap.

\subsection{Extension: models that postdate the paper}

\begin{table}[h]
\centering
\begin{tabular}{@{}lcccc@{}}
\toprule
\textbf{Model} & \textbf{parse@1} & \textbf{compile@1} & \textbf{run@1} & \textbf{run@5} \\
\midrule
deepseek-chat-v3-0324 (2025 anchor) & 63.2\% & 61.4\% & 46.3\% & 63.2\% \\
qwen3-235b-thinking-2507 & 73.7\% & 70.9\% & 52.3\% & 68.4\% \\
deepseek-v4-pro & 52.3\% & 49.5\% & 39.0\% & 61.4\% \\
\textbf{qwen3-coder} & \textbf{99.7\%} & \textbf{96.8\%} & \textbf{67.4\%} & \textbf{70.2\%} \\
\bottomrule
\end{tabular}
\caption{Same harness and 57 designs. Qwen3-Coder is 21 points above the
2025 anchor at run@1. DeepSeek-v4-pro is weaker at pass@1 because it is
the model hit hardest by the token cap (47.7\% parse-failure rate), not
because the newer checkpoint is less capable.}
\label{tab:hls-ext}
\end{table}

\begin{figure}[h]
\centering
\includegraphics[width=0.82\textwidth]{repro_vs_paper.png}
\caption{HLS-Eval pass@1 by model and metric, bootstrap 95\% CIs over
designs. Synthesis is omitted (no Vitis).}
\label{fig:hls}
\end{figure}

Qwen3-Coder saturates parse and compile, so the harness --- not the
model --- is now the bottleneck at those rungs. One year of model
progress did not translate uniformly under a fixed 8k token ceiling;
the newer, more verbose models are more exposed to it.

\subsection{What broke, independent of any model}

Found by compiling the reference designs with a stock compiler, before
any LLM call:

\begin{itemize}
\item 12 of 32 CHStone designs ship with no testbench. Runnable is
undefined for them, not failed. Scoring them as failures understates
CHStone runnability by 37.5\%.
\item \code{flowgnn/fgnn\_linear}'s testbench calls \code{printf} without
\code{<cstdio>}. Builds under Vitis's bundled clang; fails under GCC and
modern Clang.
\item \code{gnnbuilder/gather\_node\_neighbors}'s testbench writes out of
bounds. GCC's stack protector aborts; Vitis clang does not enable it.
\item The upstream harness has no execution timeout and inherits the
parent cwd, so generated infinite loops hang a worker and relative-path
testbench I/O reads the wrong directory. The Vitis-free backend used
here sets both.
\item The published leaderboard JSON in the upstream repo is still the
dummy payload. This run emits a real payload in their schema.
\end{itemize}

%----------------------------------------------------------------------
\section{Improvements I considered and did not run}
\label{sec:improvements}

The gaps above are specific enough that I could have spent another
round of API credits (and a Vitis install) chasing a closer match to
the published tables. I did not. The harness was frozen once the HLS
sweep started, and the analog cut stopped at an executable Fig.~16
rather than at a silicon-matching accuracy. What follows is the
improvement list I would actually run --- in the group, or if someone
re-runs these artifacts --- not a claim that the current numbers are
the best the models can do.

\subsection{HLS-Eval}

\begin{enumerate}
\item \textbf{Raise \code{max\_tokens} to 12k--16k and re-sweep.}
This is the single highest-leverage change. 21\% of calls in this run
ended with \code{finish\_reason=length}; DeepSeek-v4-pro's 47.7\%
parse-failure rate and the 46\% vs.\ 63\% run@1 gap versus the paper
are both consistent with truncation, not with a weaker model. A
second $n=5$ pass on the same 57 designs would likely lift pass@1
toward the paper and toward Qwen3-Coder's already-saturated parse
rate. I left it undone so the reported numbers stay tied to one
frozen harness, and so the token cap itself remains a finding rather
than something I patched away.

\item \textbf{Re-test Qwen2.5-Coder-32B on Together AI}, the paper's
actual provider. The Cloudflare/OpenRouter failure may be
provider-specific. A Together key would tell us whether the second
paper anchor is reproducible at all. I did not add that key.

\item \textbf{Pin provider and quantization} (or disable OpenRouter
load-balancing). The paper used one serving stack. This run did not.
That confound sits underneath every numeric comparison in
\S\ref{sec:hls} even with an unlimited token budget.

\item \textbf{Install Vitis HLS} and restore synthesizable, the fourth
metric. Three of four was a deliberate Vitis-free choice, not a
claim that synthesis does not matter. The stock-compiler backend
already showed that 60 reference designs need no Xilinx software for
parse/compile/run.

\item \textbf{Bring CHStone and Rosetta back} once the 12 testbench-less
CHStone kernels and Rosetta's Xilinx-header designs are handled, to
reach the paper's 94-design scope. The 57-vs-60-design gap inside the
already-clean suites should be closed in the same pass.
\end{enumerate}

None of those change the qualitative result: Qwen3-Coder is ahead of
the 2025 DeepSeek-V3 anchor under identical constraints, and the
harness --- not the student --- is now the thing to fix first.

\subsection{Analog macromodel}

The 70\% Case~I number is a digital stand-in. The obvious lifts, not
taken:

\begin{itemize}
\item Replace the integer LadderFM delay with a gm-C ladder from
TRETS ref.~[75] / the FPAA VM Xcos code.
\item Train on the full Speech Commands v2 4-word split (this run used
the mini subset) and, if licensed, compare against TIMIT/TI-digits as
the paper does.
\item Add the filter benchmark (ripple / stopband on a chirp) --- fully
specified in TRETS \S3.4, no dataset, cheap --- and Case~II 32-word.
\item Run the same Fig.~16 program on the lab's SoC FPAA, which is the
only way 70\% becomes a hardware number.
\end{itemize}

AMS-Eval (\S\ref{sec:amseval}) is the improvement I actually want to
build, not another accuracy chase on the macromodel.

%----------------------------------------------------------------------
\section{Analog benchmarks: TRETS 2024 and JLPEA 2023}
\label{sec:ams}

TRETS~\cite{trets} defines analog/mixed-signal benchmarks (acoustic,
vision, comm, filter) the way CHStone/MachSuite/Rosetta defined digital
HLS ones. JLPEA~\cite{jlpea} is the direct predecessor: same FPAA
analog-HLS flow, with a worked acoustic classifier (subbanded C$^4$ BPF,
ladder delays, VMM $+$ winner-take-all). I folded them. There is no
experimental table to match. What can be checked without an FPAA is
(i)~the scaling and CAB arithmetic, and (ii)~running Fig.~16's
level-1 program on a public dataset.

\subsection{The artifact URL was already dead at publication}

TRETS footnote~1: designs and numbers ``openly available'' at
\url{https://hasler.ece.gatech.edu/AnalogBenchmarks}. Live check
2026-08-14: 404. Wayback CDX has two captures, both 2023-11-19 (six
weeks after TRETS acceptance): HTTP 301, then HTTPS \textbf{404}. There
is no archived suite. Authority for this reproduction is the PDFs
(NSF PAR 10504737 and 10504736).

Afolabi Ige later open-sourced the \emph{synthesis} tool
(\href{https://github.com/a-folabi/ashes}{a-folabi/ashes}). That is
Python/XCOS $\to$ FPAA switch list / GDS, not the missing benchmark
files. I did not clone it.

\subsection{Scaling (Figs.~3 and 11)}

\begin{figure}[h]
\centering
\includegraphics[width=0.92\textwidth]{scaling_check.png}
\caption{TRETS Figs.~3 and 11 re-derived from stated anchors, not from
digitized pixels. Bandwidth interpolates the four caption points
exactly. CAB/FG/MAC counts scale as $1/\lambda^2$ from the 350\,nm SoC
FPAA (600k FG parameters, $\approx$1~TMAC/s, 98 CABs on a 25\,mm$^2$
reference die). If Fig.~11's drawn lines are shallower than
$\lambda^{-2}$, that is a later overlay, not a fit.}
\label{fig:scaling}
\end{figure}

Bandwidth anchors, by construction: 350\,nm $\to$ 60\,MHz, 130\,nm $\to$
500\,MHz, 40\,nm $\to$ 6\,GHz, 14\,nm $\to$ 50\,GHz. Comm Case~III
(2\,GHz on a $2\times 3$\,mm 40\,nm FPAA~\cite{trets}) is consistent
with the 6\,GHz 40\,nm anchor --- 3$\times$ margin. I did not re-read
the 40\,nm FPAA measurement paper.

\subsection{Case~II CAB arithmetic (Figs.~17--18)}

The worked example is a 12-band AmplitudeBPF, three ladder delay stages
$\to$ 48 lines, and a $48\times 128$ VMM+WTA (32 words $+$ 96 nulls)
with a 6{,}144-weight matrix. Fig.~18 costs, SoC FPAA units:

\begin{table}[h]
\centering
\begin{tabular}{@{}lrl@{}}
\toprule
\textbf{Block} & \textbf{CABs} & \textbf{Rule} \\
\midrule
AmplitudeBPF $\times$ 12 & 12 & Fig.~18: 1 CAB each \\
LadderFM, 12 ch $\times$ 3 stages & 36 & Fig.~18 LadderF rate (1 CAB/stage). LadderFM itself points at~\cite{ladder}. \\
VMMWTA $48\to 128$ & 128 & 1 WTA output / CAB on the SoC (body text) \\
\midrule
WTA only & 128 & the sentence the paper actually writes \\
Full Fig.~17 program & 176 & \\
\bottomrule
\end{tabular}
\caption{Case~II acoustic resource count. Weights $48\times 128=6144$
match the paper exactly.}
\label{tab:cab}
\end{table}

TRETS never reprints how many CABs the SoC FPAA has. Fig.~18 is
``baselined to the CABs in the SoC FPAA~\cite{socfpaa}.'' That chip has
\textbf{98 CABs}. Then 128/98 $= \mathbf{1.31}$ chips --- greater than
one, less than two. The claim ``more WTA stages than would fit on an
SoC FPAA (barely)\ldots two SoC FPAAs'' \textbf{holds}, but only if
that 98 is granted. If someone assumed $\sim$128 CABs on the SoC, WTA
would be exactly one chip and ``barely'' would not hold. The full
program is 176/98 $= 1.80$ chips, still two. The 130\,nm standard-cell
area ($250\,\mu\mathrm{m}\times 300\,\mu\mathrm{m}$) was not
re-derived.

Fig.~17 uses 12 bands. Body text elsewhere says ``16 to 32 bands.'' I
costed the figure.

\subsection{Executable Fig.~16, acoustic Case~I}

JLPEA Fig.~3 is Case~I: 16 subbands, ladder taps, VMM+WTA, 2--5 word
vocabulary. TRETS Fig.~16 is the three-line program:

\begin{verbatim}
y1  = AmplitudeBPF(In, center=f1, Q=Q1, fpeak=f2)
y2  = M2V(LadderFM(input=y1, Ibias1=Ib1, Ibias2=Ib2))
Out = VMMWTA(input=y2, W=W1, Thresh=theta1, freq=100)
\end{verbatim}

(Fig.~17's Case~II listing writes \code{freq=400}; Case~I follows
Fig.~16.) I implemented those blocks as discrete-time macromodels
(numpy/scipy; no ODEs, no FPAA VM). BPF is an RBJ biquad stand-in for
the C$^4$; AmplitudeBPF is BPF $+$ rectify $+$ envelope LPF; LadderFM
is an \textbf{integer sample delay} (25\,ms/stage) --- \code{Ibias1} and
\code{Ibias2} are accepted and ignored. That is not a gm-C ladder.
VMMWTA trains with a softmax surrogate and evaluates with hard WTA.

Dataset: TensorFlow mini Speech Commands, labels yes/no/up/down. TIMIT
is LDC-licensed and was not used. 16 log-spaced bands, 50\,Hz--4\,kHz
(Nyquist-limited relative to the paper's 10\,kHz). Case~I CAB cost is
68.9, which fits one 98-CAB SoC FPAA, consistent with ``Case~I already
demonstrated on the SoC FPAA.''

Held-out, $n=200$ test / 600 train: \textbf{accuracy 0.70} (chance 0.25;
normal-approx 95\% interval $[0.64, 0.76]$). Per class: yes 0.78, no
0.57, up 0.89, down 0.53. This is not 90\%, and it is not a silicon
number. The paper's demonstrated classifier is analog, on
TIMIT/TI-digits, with ladder dynamics this macromodel does not have.
The deliverable is that Fig.~16 runs, not that the accuracy matches
hardware.

Eight unit tests pass, including the 6{,}144-weight / 1.31-chip Case~II
arithmetic.

%----------------------------------------------------------------------
\section{A first project: AMS-Eval}
\label{sec:amseval}

HLS-Eval's contribution is a metric ladder (parseable $\to$ compilable
$\to$ runnable $\to$ synthesizable) plus a harness other people can
re-run. TRETS asks for the analog analogue and then 404s the files.
The obvious next step --- and the one I would want to do in the group
--- is to reuse that harness shape on Fig.~15's block library:

\begin{center}
\small
\begin{tabular}{@{}lll@{}}
\toprule
\textbf{HLS-Eval} & \textbf{AMS-Eval} & \textbf{Grade} \\
\midrule
parseable & parseable & AST; calls restricted to the Fig.~15 whitelist \\
compilable & simulatable & runs under the tier-F macromodel, bounded time \\
runnable & functional & accuracy / filter spec on held-out data \\
synthesizable & resource-feasible & \code{cab\_cost(program)} $\le$ Fig.~18 budget \\
\bottomrule
\end{tabular}
\end{center}

The macromodels and CAB cost function from \S\ref{sec:ams} are the
grading backend. I did not run an LLM sweep on analog programs. If
frontier models score near zero, that is itself the result (``they
cannot compose analog block programs yet''). The six HLS-Eval harness
fixes (timeout, cwd, pinned generation options, provider preflight,
retry, larger \code{max\_tokens}) would be baked in from day one rather
than discovered mid-sweep.

This sits on the Hasler--Hao intersection: her digital HLS/benchmark
methodology, his FPAA block library, an evaluation story the lab
already knows how to read.

%----------------------------------------------------------------------
\section{Papers read, not reproduced}
\label{sec:other}

\textbf{SmartPAF} (Tong, Dang, Golder, Raychowdhury, Hao, Krishna;
MLSys 2024)~\cite{smartpaf} replaces ReLU/softmax/etc.\ with tuned
low-degree polynomials for FHE inference, with claimed
$1.42$--$13.64\times$ accuracy and $6.79$--$14.9\times$ speedup on
ResNet-18 / ImageNet-1k. Public code was promised. This is the closest
\emph{numeric} analogue to HLS-Eval --- standard model, standard
dataset, concrete factors --- but it is not analog/neuromorphic and
ImageNet$+$FHE is a multi-day environment. I read it; I did not run it.
Worth doing if the group cares about private inference, not as a second
weekend project beside analog.

\textbf{CRYPTONITE}~\cite{cryptonite} auto-generates HLS accelerators
for Fiat Cryptography primitives (88.88\% resource reduction / 54.31\%
latency vs.\ naive synthesis) and is open-sourced. It wants Vitis HLS
2023.1 plus SVF/GiNaC/egg/LLVM. I went Vitis-free on HLS-Eval
specifically to avoid that dependency. Interesting as an HLS
compiler-stack paper; not AI/NN-flavored; not started.

\textbf{TODAES survey} of ML for software-assisted hardware
verification~\cite{survey} is background reading. Nothing to
reproduce.

Two IEEE Xplore links from an earlier scan did not resolve (empty
document ID; 418 on \code{11240893}).

%----------------------------------------------------------------------
\section{What I am not claiming}

\begin{itemize}
\item I did not match HLS-Eval's 94-design, 4-metric, Together AI
setup. Three metrics, 57 designs, OpenRouter, token-capped.
\item I did not run analog programs on an FPAA, and 70\% is not a
reproduction of a published accuracy. It is evidence that Fig.~16 has
executable semantics.
\item I did not digitize Fig.~11's pixels. CAB scaling is $1/\lambda^2$
from stated anchors.
\item I did not implement AMS-Eval, SmartPAF, or CRYPTONITE.
\end{itemize}

The analog work is the same genre as the HLS-Eval work: run the
benchmark definition as code, check the numbers that can be checked,
write down what broke. The proposed first project is to put an LLM in
front of the analog blocks the way HLS-Eval put one in front of C++
kernels.

%----------------------------------------------------------------------
\section*{Appendix: where the files are}

\begin{itemize}
\item HLS-Eval: \code{hls-eval/} --- report, \code{FINDINGS.md},
\code{results/leaderboard.json} (real payload, their schema), raw
per-design output, Vitis-free backend in \code{ext/stock\_tools.py}.
\item Analog: \code{analog-benchmarks/ams-bench/} ---
\code{amsblocks/}, \code{analysis/scaling.py},
\code{examples/acoustic\_case1.py}, \code{FINDINGS\_ams.md}.
\item This PDF's figures: \code{report/figures/}.
\end{itemize}

\begin{thebibliography}{9}

\bibitem{hlseval}
S.~Abi-Karam and C.~Hao,
``HLS-Eval: A Benchmark and Framework for Evaluating LLMs on High-Level
Synthesis Design Tasks,''
ICLAD 2025.
\url{https://arxiv.org/abs/2504.12268}

\bibitem{trets}
J.~Hasler and C.~Hao,
``Programmable Analog System Benchmarks Leading to Efficient Analog
Computation Synthesis,''
\emph{ACM TRETS}, vol.~17, no.~1, Art.~12, Jan.~2024.
\url{https://doi.org/10.1145/3625298}

\bibitem{jlpea}
A.~Ige, L.~Yang, H.~Yang, J.~Hasler, and C.~Hao,
``Analog System High-Level Synthesis for Energy-Efficient Reconfigurable
Computing,''
\emph{J.\ Low Power Electron.\ Appl.}, vol.~13, no.~4, p.~58, 2023.
\url{https://www.mdpi.com/2079-9268/13/4/58}

\bibitem{socfpaa}
J.~Hasler,
``A SoC FPAA with RISC microcontroller,''
\emph{IEEE JETCAS}, 2016.
TRETS ref.~[31]. 98 CABs.

\bibitem{ladder}
Ladder-filter delay paper, TRETS ref.~[75] / TCAS-I.
Not re-implemented; LadderFM is an integer delay in this run.

\bibitem{smartpaf}
J.~Tong, T.~Dang, A.~Golder, A.~Raychowdhury, C.~Hao, and T.~Krishna,
``Accurate Low-Degree Polynomial Approximation of Non-Polynomial Operators
for Fast Private Inference in Homomorphic Encryption,''
MLSys 2024.

\bibitem{cryptonite}
K.~Sharma et al.,
``CRYPTONITE: Scalable Accelerator Design for Cryptographic Primitives
and Algorithms,''
arXiv:2505.14657.
\url{https://github.com/KarthikeyaSharma16/Cryptonite}

\bibitem{survey}
Wu, Li, et al.\ (incl.\ C.~Hao),
``A Survey of Machine Learning for Software-Assisted Hardware Design
Verification,''
\emph{ACM TODAES}, 10.1145/3661308.

\end{thebibliography}

\end{document}
